% fairscope software paper -- DRAFT for review, not submitted.
% arXiv categories: cs.LG (primary), cs.CY (cross-list).
% Compile: pdflatex paper.tex
\documentclass[11pt]{article}
\usepackage[utf8]{inputenc}
\usepackage[T1]{fontenc}
\usepackage[margin=1in]{geometry}
\usepackage{amsmath}
\usepackage{booktabs}
\usepackage{url}
\usepackage[hidelinks]{hyperref}

\title{\texttt{fairscope}: subgroup-stratified, calibration-aware fairness\\
auditing for machine-learning models}
\author{Rajveer Singh Pall\thanks{ORCID: \url{https://orcid.org/0009-0001-6762-6134}}}
\date{June 27, 2026}

\begin{document}
\maketitle

\begin{abstract}
Fairness audits in applied machine learning repeatedly require statistical machinery that
mainstream toolkits do not expose as first-class, subgroup-stratified functions: confidence
intervals for per-subgroup AUC, per-subgroup calibration error, and significance tests for
differences in subgroup performance. We present \texttt{fairscope}, an open-source,
test-driven Python library that packages these peer-reviewed methods behind a small, typed
API, and introduces one novel module---a five-axis \emph{Cross-Platform Fairness Evaluation}
(CPFE) protocol that jointly characterises discriminative, calibration,
statistical-significance, prediction-equity, and attribution-stability failures when a model
is deployed on a distribution different from the one it was trained on. \texttt{fairscope} is
an \emph{engineering} contribution grounded in prior published work: every statistical
primitive ports a documented method and cites its source; only the CPFE protocol is presented
as novel. The library ships with 100\% line coverage on its statistical core, continuous
integration across Python 3.9--3.12, and small committed fixtures that regression-test the
\emph{direction and approximate magnitude} of published findings. It is available on PyPI
(\texttt{pip install fairscope}) under the MIT license.
\end{abstract}

\section{The gap}

Two widely used fairness toolkits anchor current practice. AIF360~\cite{aif360} provides a
broad set of group-fairness metrics and bias-mitigation algorithms; Fairlearn~\cite{fairlearn}
focuses on constraint-based mitigation and disparity metrics. Both are valuable, and
\texttt{fairscope} is complementary rather than competitive. However, recurring needs in
applied audits are not available in either as first-class, subgroup-stratified functions:

\begin{itemize}
  \item \textbf{Uncertainty on subgroup discrimination.} Reporting a per-subgroup AUC without
    a confidence interval invites over-interpretation of noise. The DeLong
    estimator~\cite{delong} gives an analytic AUC variance, but neither toolkit exposes a
    per-subgroup DeLong interval as a primary function.
  \item \textbf{Calibration parity.} Two subgroups can have equal accuracy yet unequal
    calibration. Per-subgroup Expected Calibration Error (ECE) is not a first-class output of
    either toolkit.
  \item \textbf{Significance of subgroup gaps.} Deciding whether an observed subgroup gap is
    distinguishable from sampling noise requires a paired or unpaired test with
    multiple-comparison control; this is left to the user.
\end{itemize}

A 2025 toolbox, \emph{meval}~\cite{meval}, moves in a similar direction by providing
fine-grained, stratified performance analysis, and we cite it as the most closely related
recent work; the precise overlap should be assessed against its current release.
\texttt{fairscope}'s distinguishing surface is the combination of subgroup DeLong intervals,
per-subgroup calibration and \emph{recalibration}, multiple-comparison-corrected significance,
the five-axis cross-platform protocol, and a per-node federated audit, packaged as a single
tested library.

\emph{Verification note: the capability comparison in Section~\ref{sec:compare} is to be
re-confirmed against the installed versions of AIF360 and Fairlearn, and meval's documented
coverage, before submission. Claims that no longer hold will be softened or removed.}

\section{Design}

\texttt{fairscope} is organised as a small set of focused modules built on NumPy, SciPy, and
scikit-learn, with optional heavy dependencies isolated behind extras: \texttt{core/}
(\texttt{delong}, \texttt{bootstrap}, \texttt{calibration} with recalibration,
\texttt{correction}, \texttt{metrics}); \texttt{healthcare/} (a one-call clinical fairness
audit and report); \texttt{nlp/} (the CPFE protocol and gradient-saliency attribution
stability); and the planned \texttt{lending/} and \texttt{federated/} modules.

\paragraph{Statistical primitives, each citing its source.} \texttt{delong} implements AUC
confidence intervals and paired/unpaired tests using the $O(n \log n)$ midrank algorithm of
Sun and Xu~\cite{sunxu} for the DeLong~\cite{delong} variance; \texttt{bootstrap} provides a
stratified bootstrap AUC-difference test; \texttt{calibration} computes ECE, maximum
calibration error, and reliability diagrams, and exposes a subgroup-stratified \emph{interface}
to two standard recalibration methods---temperature scaling~\cite{guo} and isotonic
regression~\cite{zadrozny} (the methods are standard; the contribution is the per-subgroup
interface and pre/post-ECE reporting); \texttt{correction} implements Bonferroni and
Benjamini--Hochberg~\cite{bh} corrections, validated to match \texttt{statsmodels};
\texttt{metrics} provides subgroup AUC/Brier/F1, symmetric disparate impact, and
equalized-odds difference.

\paragraph{Engineering discipline.} Every public function has full type hints, a NumPy-style
docstring with a runnable example, explicit input validation (no silent failures: an AUC on a
single-class subgroup raises rather than returning a meaningless value), and a fixed
\texttt{random\_state} wherever results are stochastic. Regression tests are the project's
credibility anchor: where an authoritative reference value exists (DeLong's worked example;
\texttt{statsmodels} multiple-testing routines), tests assert against it. The statistical core
has 100\% line coverage; CI runs pytest, coverage, ruff, and black across Python 3.9--3.12.
Heavy NLP dependencies (torch, transformers, captum) install only via \texttt{pip install
fairscope[nlp]}; SHAP via \texttt{fairscope[shap]}. Datasets are not bundled.

\section{The CPFE protocol}

The one novel module is the \textbf{Cross-Platform Fairness Evaluation (CPFE)} protocol.
Models trained on one data source and deployed on another can degrade in ways a single AUC
number cannot distinguish. CPFE evaluates five orthogonal axes, treating each non-training
platform against a reference:

\begin{enumerate}
  \item \textbf{Discriminative performance}---macro one-vs-rest AUC and macro F1, within- and
    cross-platform, with the relative change $\Delta$AUC\%.
  \item \textbf{Calibration}---confidence--accuracy ECE per platform (the standard multiclass
    ECE of Guo et al.~\cite{guo}).
  \item \textbf{Statistical significance}---a stratified \emph{bootstrap} standard error on the
    macro-AUC difference between platforms (independent test sets), with Bonferroni correction
    across platform pairs.
  \item \textbf{Prediction equity}---symmetric disparate impact and equalized-odds difference
    per class, treating platform membership as the group.
  \item \textbf{Attribution stability}---Jaccard overlap of the top-$K$ gradient-saliency
    token sets across platforms, $s_i = \|\partial P(y\,|\,x)/\partial E_i\|_2$,
    $K \in \{5,10,15,20\}$.
\end{enumerate}

A \texttt{CPFEReport.deployment\_readiness()} method returns a structured per-axis screening
diagnostic. It is explicitly a diagnostic, not a deployment decision or compliance verdict.
Each axis reports its value, the threshold applied, and the threshold's \emph{source}: the
disparate-impact four-fifths rule and the calibration bands are reference values stated in the
underlying study; the discrimination $\Delta$AUC limit is an \emph{illustrative,
user-configurable} parameter, because that study explicitly declines to set a validated
$\Delta$AUC decision threshold.

\section{Worked replications}

We report two demonstrations. The first reproduces a \emph{published} finding's direction and
magnitude; the second exercises the CPFE protocol on a synthetic shift. \textbf{No real
restricted datasets or trained models are distributed; committed fixtures are small and
explicitly labelled.}

\paragraph{4.1 Healthcare audit (direction + magnitude of a published result).} The author's
published diabetes study~\cite{diabetes} reports, on an external validation cohort (BRFSS,
$n=1{,}285{,}783$), a substantially lower AUC for elderly adults ($\geq 60$) than for younger
adults---$0.607$ versus $0.742$, a gap of $0.135$ with $p < 0.001$. \texttt{fairscope}'s
\texttt{HealthcareFairnessAudit} is regression-tested on a \emph{synthetic fixture calibrated
to reproduce that direction and approximate magnitude} (it is not real BRFSS data and does not
reproduce the full run): the audit recovers a young-vs-elderly AUC gap of $0.133$ (young
$0.764$, elderly $0.631$) with a Bonferroni-adjusted $p < 0.01$. The test asserts the gap
within a tolerance band and a conservative significance threshold, and documents that the
published $p < 0.001$ was obtained on the full $1.28$M-row cohort, not the fixture.

\paragraph{4.2 CPFE protocol (synthetic cross-platform shift).} On a synthetic two-platform
fixture---a high-signal reference platform and a platform with an injected
over-confident-but-wrong shift---the protocol's five axes behave as designed (these are the
library's actual outputs on the fixture, reproducible via \texttt{paper/cpfe\_demo.py}, seed
$20260627$, not real-world measurements): macro-AUC falls from $0.997$ to $0.341$
($\Delta$AUC $-65.8\%$); macro-F1 falls from $0.958$ to $0.019$; ECE rises from $0.172$ to
$0.776$; the bootstrap macro-AUC difference is significant (Bonferroni-adjusted $p \ll 0.001$,
$z = 163.2$); and a synthetic attribution comparison with near-disjoint token sets yields
$J@5 = 0.11$. Notably, the prediction-equity axis shows \emph{no} four-fifths violation on this
fixture (per-class disparate impact $0.97$--$0.99$), because the degraded predictions remain
roughly proportional across classes---illustrating that the five axes capture genuinely
distinct failure modes, and that a single discrimination metric would have masked which
dimensions failed.

\section{Comparison with existing toolkits}
\label{sec:compare}

\emph{Provisional---to be re-verified against installed package versions before submission, per
Section~1.}

\begin{table}[h]
\centering
\begin{tabular}{lccc}
\toprule
Capability & AIF360 & Fairlearn & \texttt{fairscope} \\
\midrule
Per-subgroup AUC with DeLong CI & not first-class & not first-class & yes \\
Per-subgroup Expected Calibration Error & not first-class & not first-class & yes \\
Subgroup significance test (+ correction) & not first-class & not first-class & yes \\
Subgroup-stratified recalibration interface & partial & partial & yes \\
Cross-platform five-axis protocol (CPFE) & no & no & yes (novel) \\
Per-node / federated audit & no & no & planned \\
Bias-mitigation algorithms & yes & yes & out of scope \\
\bottomrule
\end{tabular}
\caption{Provisional capability comparison (see verification note).}
\end{table}

\texttt{fairscope} does not implement bias-mitigation or constraint-based training; for those,
AIF360 and Fairlearn remain the appropriate tools. \texttt{fairscope}'s scope is rigorous,
uncertainty-aware \emph{measurement}.

\section{Limitations}

\begin{itemize}
  \item The committed healthcare fixture is \textbf{synthetic}, calibrated to a published
    direction and magnitude; it is a pipeline regression test, not an empirical result.
  \item CPFE's per-axis screening thresholds are descriptive: the disparate-impact and
    calibration bands are reference values from a single study, and the $\Delta$AUC limit is
    illustrative and user-configurable, not a validated decision threshold. None should be
    treated as regulatory criteria.
  \item Gradient-saliency attribution is a local linear approximation and can be unstable
    across seeds; the Jaccard axis should be read as consistent-with-instability, not as
    definitive.
  \item The capability comparison (Section~\ref{sec:compare}) requires confirmation against
    current toolkit releases before publication.
\end{itemize}

\section{Availability}

Source: \url{https://github.com/Rajveer-code/fairscope} (MIT license). Install: \texttt{pip
install fairscope}; NLP extras via \texttt{pip install fairscope[nlp]}. The release is archived
and citable via \texttt{CITATION.cff}. Continuous integration, 100\% statistical-core coverage,
and a public, incremental commit history accompany the code.

\paragraph{Note on the CPFE source study.} The five-axis design is the author's; the
cross-platform mental-health NLP study that motivates it is under review, and specific
unpublished figures from it are intentionally not cited here. The citation will be added once
that paper is public.

\begin{thebibliography}{9}

\bibitem{aif360} Bellamy, R.~K.~E., et al. (2019). AI Fairness 360: An extensible toolkit for
detecting and mitigating algorithmic bias. \emph{IBM Journal of Research and Development},
63(4/5). arXiv:1810.01943.

\bibitem{bh} Benjamini, Y., \& Hochberg, Y. (1995). Controlling the false discovery rate.
\emph{Journal of the Royal Statistical Society: Series B}, 57(1), 289--300.

\bibitem{fairlearn} Bird, S., et al. (2020). Fairlearn: A toolkit for assessing and improving
fairness in AI. \emph{Microsoft Technical Report} MSR-TR-2020-32.

\bibitem{delong} DeLong, E.~R., DeLong, D.~M., \& Clarke-Pearson, D.~L. (1988). Comparing the
areas under two or more correlated receiver operating characteristic curves: a nonparametric
approach. \emph{Biometrics}, 44(3), 837--845.

\bibitem{guo} Guo, C., Pleiss, G., Sun, Y., \& Weinberger, K.~Q. (2017). On calibration of
modern neural networks. \emph{ICML}, PMLR 70:1321--1330.

\bibitem{meval} Sutariya, D., \& Petersen, E. (2025). meval: A Statistical Toolbox for
Fine-Grained Model Performance Analysis. arXiv:2512.17409.

\bibitem{diabetes} Pall, R.~S., Yadav, S., Bhalerao, S., Sahu, S., Ahluwalia, R., \&
Awadhiya, B. (2026). Comprehensive Evaluation of Machine Learning for Type~2 Diabetes Risk
Prediction: Large-Scale External Validation and Fairness Analysis. \emph{IEEE CIPHER 2026}.
DOI: 10.1109/CIPHER70417.2026.11523789.

\bibitem{sunxu} Sun, X., \& Xu, W. (2014). Fast implementation of DeLong's algorithm for
comparing the areas under correlated receiver operating characteristic curves. \emph{IEEE
Signal Processing Letters}, 21(11), 1389--1393.

\bibitem{zadrozny} Zadrozny, B., \& Elkan, C. (2002). Transforming classifier scores into
accurate multiclass probability estimates. \emph{KDD}, 694--699.

\end{thebibliography}

\end{document}
